\documentclass{article}
\usepackage{amsmath}

\title{Ch. 11 CW Solutions}
\author{Principles of Microeconomics}
\date{\today}

\begin{document}

\maketitle

\begin{enumerate}

\item

	\begin{enumerate}

	\item It is. Once a movie is on the TV, it's difficult to prevent people in the room who didn't pay for the movie from watching, so a movie showing is not excludable. Additionally, one person watching the movie doesn't diminish another person's ability to watch the movie, so a movie showing is not rival. 

	\item They should stream three movies. Their collective willingness to pay drops below \$8 after the third movie.
	
	\item They would spend $\$8 \cdot 3 = \$24$ total, so each person would spend \$6. Therefore, each person's consumer surplus would be:
	
		\begin{align*}
		\text{Dwayne: } &CS = 7 + 6 + 5 - 6 = 12 \\
		\text{Javier: } &CS = 5 + 4 + 3 -6 = 6 \\
		\text{Salman: } &CS = 3 + 2 + 1 - 6 = 0 \\
		\text{Chris: } &CS = 2 + 1 - 6 = -3
		\end{align*}
		
	\item Everyone could pay the same fraction of their willingness to pay The difficulty is getting people to tell the truth about their willingness to pay if they know it will determine what share of the cost they'll bear.
	
	\item He will not. He'll have an incentive to overstate his willingness to pay in an effort to get them to watch more movies. Each additional movie only costs him \$2 because the roommates split the \$8 cost four ways. His willingness to pay for a movie is greater than \$2 for a fourth and fifth movie, so he wants them watch at least five movies.
	
	\item It's difficult to find the optimal number of movies and then pay for them in a fair way.

	\end{enumerate}
	
\item

	\begin{enumerate}
	
	\item Basic research is a public good because it's not excludable or rival. Once it's published, you can prevent people who didn't fund the research from using it. Additionally, one person's use of the research doesn't diminish another person's ability to use it.
	
	\item The United States funds a lot of research at universities which it pays for using tax revenue.
	
	\item It's not totally consistent with my argument that basic research is not excludable. You could argue that there are barriers to foreign firms using domestic research. For one, if the research is in English, it will take time to translate the findings into other languages. For two, firms that are geographically closer to the researchers might have more exposure to their findings and an easier time adopting the advances. As a result, basic research might be slightly excludable. That said, over time the research will desseminate, and it will become non-excludable. 
	
	\end{enumerate}
	
\item

	\begin{enumerate}
	
	\item They would. The collective willingness to pay for the fireworks exceeds the cost.
	
	\item If they split the cost equally, everyone would pay \$120 for the fireworks. In that case, Callie would vote in favor because her willingness to pay is greater than \$120, but everyone else would vote no because their willingness pay is less than \$120. As a result, the vote would not yield the same answer as the cost-benefit analysis.
	
	\item No, the collective willingness to pay for the fireworks is less than the cost.
	
	\item Again, each person would pay \$120. Ned and Bess would both vote in favor because their willingness to pay is greater than \$120. Nancy would vote against because her willingness to pay is less than \$120. As a result, the vote would not yield the same answer as the cost-benefit analysis.
	
	\item It's difficult to decide whether or not to supply them. 
	
	\end{enumerate}
	
\item During rush hour, the transportation system becomes overcrowded, so not everyone can fit on the trains. That means the trains become rival, and one person's use diminishes another person's ability to use them. In other words, the trains become a common resource that get overused and create a negative externality. By increasing the price of using the transportation system during rush hour, the government is reducing the free-market quantity of train usage down towards its efficient quantity.

\item

	\begin{enumerate}
	
	\item \begin{align*}
		Q_{D,1} &= 5 - \frac{1}{2}P \\
		\frac{1}{2}P_1 &= 5 - Q \\
		P_1 &= 10 - 2Q \\
		Q_{D,2} &= 3 - \frac{1}{2}P \\
		\frac{1}{2}P_2 &= 3 - Q \\
		P_2 &= 6 - 2Q \\
		P_{mkt} &= P_1 + 3P_2 \\
		&= 10 - 2Q + 3(6 - 2Q) \\
		&= 28 - 8Q \\
		20 &= 28 - 8Q_{opt} \\
		Q_{opt} &= 1
		\end{align*}
	
	\item Zero, they're max willingnesses to pay are \$10 and \$6 from their inverse demand equations. That's less than the price of a single hour.
	
	\end{enumerate}

\end{enumerate}

\end{document}